\section{Einstein linear e normalização histórica da massa}
\label{sec:fund-normalizacao}

A reprodução geométrica depende de \(\omega\), \(k\) e dos vínculos.
Relacioná-la a uma massa em unidades físicas exige ainda conferir o
coeficiente da equação de propagação. Nesta seção, \(\Box\) designa o
d'Alembertiano de segunda ordem. O símbolo \(\Box^2\) utilizado no TG
representa esse operador, não sua composição de quarta ordem.

Uma verificação independente, com dez componentes métricas e quatro
componentes de \(q_\mu\) livres, parte da conexão e contrai Ricci e
Einstein. Definindo \(B_\nu=q^\alpha\bar H_{\alpha\nu}\), obtém-se
\begin{equation}
 G_{\mu\nu}^{(1)}=-\frac12\Box\bar H_{\mu\nu}
 -\frac12(q_\mu B_\nu+q_\nu B_\mu)
 +\frac12\eta_{\mu\nu}q^\alpha B_\alpha,
 \label{eq:fund-einstein-geral}
\end{equation}
com \(\Box\mapsto-q^\alpha q_\alpha\) na amplitude de Fourier.
A contração \(q^\mu G_{\mu\nu}^{(1)}\) é identicamente nula.
Quando \(B_\nu=0\), resulta
\(G_{\mu\nu}^{(1)}=-\Box\bar H_{\mu\nu}/2\). A convenção oposta
de Riemann usada nas equações históricas inverte esse sinal:
\begin{equation}
 G_{\mu\nu}^{(1),\mathrm{hist}}=
       \frac12\Box\bar H_{\mu\nu}.
 \label{eq:fund-einstein-historico}
\end{equation}
A identidade foi conferida para as duas assinaturas de Minkowski.
Mudar a assinatura modifica a expressão de \(\Box\), mas não elimina
o fator \(1/2\).

As páginas impressas 23--24 do TG fornecem, nas equações (4.8) e
(4.9), o coeficiente linear
\(T_{\mu\nu}^{\mathrm{mass}}
=-m_{\mathrm{TG}}^2c^2\bar H_{\mu\nu}/(8\pi G\hbar^2)\)
e \(G_{\mu\nu}=8\pi GT_{\mu\nu}^{\mathrm{mass}}/c^4\) em vácuo
\autocite{tg2003}. Sua substituição literal em
\eqref{eq:fund-einstein-historico} produz
\begin{equation}
 \left(\Box+\frac{2m_{\mathrm{TG}}^2}{\hbar^2c^2}\right)
 \bar H_{\mu\nu}=0,
 \label{eq:fund-tg-literal}
\end{equation}
reproduzindo o fator dois das equações (4.17)--(4.19). Há, entretanto,
uma questão dimensional: se \(m_{\mathrm{TG}}\) é massa SI e
\(x^0=ct\), então
\begin{equation}
 \left[\frac{m_{\mathrm{TG}}^2}{\hbar^2c^2}\right]=L^{-6}T^4,
 \qquad
 \left[\left(\frac{m_gc}{\hbar}\right)^2\right]=L^{-2}.
 \label{eq:fund-dimensoes-massa}
\end{equation}
Falta um fator \(c^4\) no primeiro coeficiente relativamente à dimensão
exigida. Reinterpretar o símbolo como massa previamente reescalada
exigiria explicitar essa definição e traduzir também a velocidade
impressa após (4.22). Usar \(c=1\) oculta a questão dimensional,
sem demonstrar a equivalência das identificações de massa.

O artigo de 2004 adota assinatura \((-+++)\) e escreve \(c^6\) no
numerador de \(T^{\mathrm{mass}}\), nas equações (4)--(5), restaurando
a dimensão de densidade de energia \autocite{depaula2004}. Sua equação
(6) usa \(G=-8\pi GT^{\mathrm{mass}}/c^4\). Denotando por
\(m_a\) o parâmetro dessa fonte e conservando o sinal de Ricci da
equação (10), a cadeia literal resulta em
\begin{equation}
 \left(\Box-2\frac{m_a^2c^2}{\hbar^2}\right)\bar H_{\mu\nu}=0.
 \label{eq:fund-artigo-literal}
\end{equation}
Já a equação (12) imprime o coeficiente \(-m_a^2c^2/\hbar^2\).
Portanto, as expressões consultadas apresentam uma discrepância de
fator dois se o parâmetro de massa for o mesmo. A mudança de assinatura
explica o sinal relativo ao TG, mas não esse fator. Reduzir o
coeficiente fonte pela metade ou definir
\(m_{\mathrm{disp}}^2=2m_a^2\) são escolhas de normalização distintas;
não se devem aplicar ambas como compensações silenciosas.

Este é um diagnóstico algébrico das equações consultadas, não uma
errata editorial confirmada ou uma reconstrução da ação não linear.
A dissertação adota a \emph{massa de dispersão} \(m_g\), definida por
\eqref{eq:fund-dispersao-si}, e \(\mu=m_gc/\hbar\).
Esse coeficiente operacional também aparece na formulação linear de
Visser \autocite{visser1998}. Sua identificação com um parâmetro de
outra normalização deve ser demonstrada em cada comparação. As tabelas
em \(\mathrm{eV}/c^2\) e a resposta futura de PTA usarão a definição
operacional; a reprodução das expressões de curvatura permanecerá em
\(\omega\) e \(k\), preservando as fontes históricas.
